\section{Google、OpenAI 与 Search：入口权如何重分配}

上一章说明 Google 和 OpenAI 都在核心牌桌上，本章专门讨论两者的入口竞争。节目中最有张力的判断是：Google 王者归来不等于危机解除，ChatGPT 也不只是一个边缘聊天工具。二者最终可能在搜索、个人助理、移动端和广告预算上直接相遇。

\subsection{Gemini 崛起后，OpenAI 为什么仍然强}

广密对 OpenAI 的看法偏乐观，原因有三。第一，OpenAI 可能是过去二十年里少数真正有机会挑战 Google 的公司。微软 Bing 也没有从根本上改变搜索入口，但 ChatGPT 把用户心智从“搜索答案”迁移到“让 AI 帮我理解、生成、执行”。第二，OpenAI 仍是最好的 AI 标的之一：人才密度、产品感、品牌和用户规模都很强。第三，OpenAI 仍有较高概率参与下一范式探索，这对长期估值很关键。

Google 的优势同样真实：Gemini、多模态、Android、Search、YouTube、Workspace、Cloud、TPU、Waymo 等资产形成了很厚的底盘。问题在于，Google 的强并不自动消灭 OpenAI；这可能是一个双巨头甚至多巨头长期并存的市场。

\subsection{Search 会怎样被 Agent 重写}

前面说 Google 的危机没有真正解除，问题就落到 Search 上：搜索不是一个单独功能，而是广告、网页分发、商家流量和用户决策的总入口。本节把传统搜索流程和 Agent 流程并排看，才能看出流量重分配为什么会触及 Google 的利润核心。

搜索的传统流程是：用户输入关键词，搜索引擎返回链接、摘要和广告，用户再点击、比较、购买或继续搜索。Agent 流程则可能是：用户给目标，Agent 搜索、比较、调用工具、议价、下单、支付、跟踪，最后返回结果。两者差异不只是界面，而是价值链位置。

\begin{figure}[H]
\centering
\includegraphics[width=0.96\textwidth]{search-agent-disruption.png}
\caption{Search 与 Agent 的流量重分配：从搜索链接到任务完成。自制概念图，依据 00:31:20--00:35:08 与 00:50:57--00:52:53 对谈内容整理。}
\end{figure}

\begin{knowledgebox}{读图：Agent 抢的是分配权}
图中用户目标先进入 Agent，而不是直接进入搜索框或 Super App。Agent 再决定调用搜索、票务、出行、支付或企业系统。谁控制这个入口，谁就有机会重新分配广告、交易、抽成和用户数据。
\end{knowledgebox}

\begin{warningbox}{Agent 不会瞬间消灭搜索}
Agent 要真正改写搜索广告，必须解决引用可信度、实时信息、支付、权限、商家接入、反作弊、审计和责任归属。短期更可能是混合形态：传统搜索融合 AI，ChatGPT 融合搜索，二者互相侵蚀。
\end{warningbox}

\subsection{ChatGPT 的移动端心智与 Gemini 的生态心智}

接下来需要从入口机制转到用户心智，因为同样能回答问题的模型，最终可能在完全不同的时间、设备和场景里被打开。本节讨论的是使用频率和生活位置，而不只是模型能力。

节目中提到一个非常产品化的观察：ChatGPT 更像生活助理，很多用户会在周末、移动端和日常生活中使用；Gemini 在一些场景中更像工作时间的生产力工具，且依托 Android 和新兴市场有天然分发优势。这个差异意味着，模型能力之外，使用频率、场景温度和用户心智同样重要。

如果 ChatGPT 能成为 10--20 亿日活级别的个人助理，它的商业化根基会非常稳；如果 Gemini 能把 Android、Search、Workspace 和多模态真正串起来，Google 也可能在 AI 时代继续保持强势。两条路径不是互斥，而是在争夺用户最先打开哪个入口。

\subsection{本章小结}

Google 的危机从“会不会错过 AI”变成“能不能在 Agent 入口里继续掌握分配权”。OpenAI 的挑战从“能不能做出好模型”变成“能不能把 ChatGPT 扩展成高频、可靠、可变现的个人助理和搜索/执行入口”。这一章的结论是：入口战争刚开始，模型榜单只是其中一小部分。

